\centerline{\bf \Huge Комбинаторный разнобой 2}

\problem{1}
В стране $2 n+1$ город. Известно, что из каждого города кроме столицы выходят по две дороги. Также известно, что из каждого города можно доехать до столицы. Сколько дорог может выходить из столицы? Движение по всем дорогам двустороннее, и любые два города соединены не более, чем одной дорогой.

\problem{2}
Три жулика, каждый с двумя чемоданами, находятся на одном берегу реки, через которую они хотят переправиться. Есть трёхместная лодка, каждое место в ней может быть занято либо человеком, либо чемоданом. Никто из жуликов не доверит свой чемодан спутникам в своё отсутствие, но готов оставить чемоданы на безлюдном берегу. Смогут ли они переправиться?

\problem{3}
В вершинах куба как-то расставлены 3 единички и 5 нулей. Каждый ход можно выбрать ребро и добавить к вершинам на его концах по 1 . Можно ли такими операциями добиться, чтобы все числа в вершинах куба делились на 7 ?

\problem{4}
Докажите, что существует граф с $2 n$ вершинами, степени которых равны $1,1,2,2, \ldots, n, n$.

\problem{5}
По кругу стоят 100 гномов разного роста. Время от времени один из них перебегает на другое место (между какими-то двумя гномами). Гномы хотят встать по росту в порядке возрастания по часовой стрелке. Какого наименьшего количества таких перебежек им заведомо хватит, как бы они не стояли изначально?

\problem{6}
Каждая точка плоскости покрашена в один из 2022 цветов. Докажите, что найдётся прямоугольник с 4 одноцветными вершинами.

\problem{7}
На клетчатой бумаге нарисован выпуклый пятиугольник. Все его вершины расположены в узлах сетки. Докажите, что найдется еще один узел сетки, расположенный (а) внутри пятиугольника или на его границе; (б) строго внутри пятиугольника.

\problem{8}
В волшебной стране несколько городов, некоторые из которых соединены дорогой. Саруман и Гэндальф борются между собой за господство в этой стране. В первый день каждому из них позволяется захватить по городу (Саруман захватывает первый). Далее, каждый день волшебникам разрешается по очереди (начинает Саруман) захватить по городу так, чтобы вновь захваченный город был соединён с уже захваченными ими городами дорогой. Запрещается захватывать города, ранее захваченные другим волшебником. Процесс заканчивается, когда все города кем-то захвачены. Существует ли такое расположение городов и дорог, что Гэндальф сможет захватить более $66 \%$ городов, вне зависимости от действий Сарумана?

\problem{9}
Можно ли шахматным конём обойти доску размером 4 на 2022, побывав в каждой клетке ровно один раз?